\documentclass[12pt]{article}

% Match the earlier reports using locally available fonts.
\usepackage[T1]{fontenc}
\usepackage{times}
\usepackage{microtype}
\usepackage{amsmath}
\usepackage[parfill]{parskip}
\setlength{\parindent}{0pt}
\setlength{\parskip}{2ex plus 0.2ex minus 0.2ex}
\usepackage[letterpaper,margin=1in]{geometry}
\usepackage[hidelinks]{hyperref}

\begin{document}

\begin{flushleft}
\textbf{Week 10: Deep Probabilistic Models} \\
Jaisel Singh (js6897) \\
November 18, 2026
\end{flushleft}

\vspace{0.1in}

What stood out to me in Kingma and Welling was that the encoder learns how to
perform approximate inference across observations. The generative model
$p_\theta(z)p_\theta(x\mid z)$ describes how latent variables produce data,
while $q_\phi(z\mid x)$ approximates the posterior in the reverse direction.
Sharing the encoder's parameters across examples replaces a separate
optimization of variational parameters for every new observation with a learned
mapping. I found this connection to week 6 especially useful: the encoder
returns a distribution over possible codes, and the decoder specifies a
distribution over observations. Their roles remain distinct even though both
are trained together using the evidence lower bound.

The reparameterization used by Kingma and Welling and by Rezende et al. made the
computation clearer to me. For a diagonal Gaussian encoder, we can write
\mbox{$z=\mu_\phi(x)+\sigma_\phi(x)\odot\epsilon$}, with
\mbox{$\epsilon\sim\mathcal{N}(0,I)$}. The randomness comes from a distribution that
does not depend on the encoder parameters, allowing gradients to pass through
the transformation for a sampled noise value. This connects stochastic
optimization to probabilistic inference without requiring analytic expectations
for the whole network. Ranganath et al.'s deep exponential families also show
that layers of random variables can use distributions beyond the Gaussian,
connecting these deeper representations to week 8's modeling choices.

Wang, Blei, and Cunningham complicated my sense of what successful training
means. In exact posterior collapse, $p_\theta(z\mid x)=p(z)$: observing the data
does not update the latent variable. They connect this to a likelihood that is
constant in the latent variable at the fitted parameters, so the problem can
exist even with exact inference. They also distinguish this from collapse
caused by an overly restrictive variational approximation. I find that
distinction important because improving the encoder cannot recover information
that the fitted generative model never uses. My remaining question is how to
evaluate whether a learned representation contains useful information about
the data beyond obtaining convincing samples or a good reconstruction score.

\vspace{0.05in}
{\footnotesize
\textit{Readings:} Kingma and Welling (2013), \href{https://arxiv.org/abs/1312.6114}{\emph{Auto-Encoding Variational Bayes}};
Rezende et al.\ (2014), \href{https://www.cs.columbia.edu/~blei/fogm/readings/RezendeMohamedWierstra2014.pdf}{\emph{Stochastic Back-propagation and Variational Inference in Deep Latent Gaussian Models}};
Ranganath et al.\ (2015), \href{https://arxiv.org/abs/1411.2581}{\emph{Deep Exponential Families}};
Wang et al.\ (2021), \href{https://www.cs.columbia.edu/~blei/papers/WangBleiCunningham2021.pdf}{\emph{Posterior Collapse and Latent Variable Non-identifiability}}.
\par}

\end{document}

%%% Local Variables:
%%% mode: latex
%%% TeX-master: t
%%% End:
